\documentclass[11pt]{article}
\usepackage[margin=1in]{geometry}
\usepackage[T1]{fontenc}
\usepackage[utf8]{inputenc}
\usepackage{amsmath}
\usepackage{amssymb}
\usepackage{amsthm}
\usepackage{booktabs}
\usepackage{enumitem}
\usepackage{hyperref}
\usepackage{url}

\newtheorem{definition}{Definition}
\newtheorem{lemma}{Lemma}
\newtheorem{theorem}{Theorem}
\newtheorem{claim}{Claim}

\title{Agentic Synthesis against Counterexample-Supplemented Sketches}
\author{Muness Castle \and Eric Rubeck}
\date{\today}

\begin{document}
\maketitle

\begin{abstract}
Coding agents are useful synthesizers but weak custodians of intent: a prompt can describe
what to build while leaving unstated which plausible implementation would be wrong.  We
present \emph{agentic synthesis against counterexample-supplemented sketches}, a method
for programming with agents under checkable rules.  A human writes a partial program---a
\emph{sketch}---that fixes strategy and names holes; a finite corpus $E$ stores observations
and promoted counterexamples; an agent edits the sketch, deterministic code, and prompts;
and a gate refuses completion until replay and semantic comparison pass every case in $E$.

The method imports the loop grammar of sketching and counterexample-guided inductive
synthesis (CEGIS) while changing the synthesizer.  In classical Sketch, a solver fills holes
inside a formal language.  Here, the synthesizer is an IDE-shaped coding agent whose edits
span ordinary source code and LLM prompts.  The result is not a proof of correctness over all
inputs.  It is a finite-corpus correctness result: if the gate passes, the current artifact is
correct for the promoted corpus under the replay and comparison semantics encoded by the
repository.  We state this guarantee explicitly, prove it from the gate definition, and show
how it is useful despite being bounded.

The running example is a clean-room RosterSynth slice: a roster remediation task where a
naive append closes the hours math but violates duplicate-booking policy.  The
counterexample promotes a sketch clause, drives an Oracle~A implementation, aligns an
Oracle~B prompt/cassette path, and is checked by replay plus compare.  The example is not
the contribution; it is the reference artifact readers can run, inspect, and adapt.
\end{abstract}

% ---------------------------------------------------------------------------
\section{Introduction}
\label{sec:introduction}

A coding agent can produce a working patch from a vague request.  That is both the appeal
and the risk.  The agent is optimized to make progress under underspecification; the human
is then left to decide whether the patch captured the rule, merely satisfied the visible test,
or invented an accidental convention.  The dangerous failure is not a syntax error.  It is a
plausible implementation that closes the obvious gap while violating the policy the prompt
failed to make explicit.

This paper describes a discipline for using agents as synthesizers without treating chat
history as the specification.  The discipline has three visible artifacts:

\begin{enumerate}[leftmargin=*]
  \item a \textbf{sketch} $\mathcal{S}$: a partial program in prose and code-shaped structure;
  \item a \textbf{counterexample-supplemented corpus} $E$: examples and promoted failures;
  \item a \textbf{gate}: executable validation over all of $E$.
\end{enumerate}

The agent edits ordinary repository files: the sketch, deterministic implementation code,
and prompts for LLM-mediated holes.  The gate is the boundary of done.  When the gate
finds a failure, the failure is not treated as a one-off bug report.  It is promoted into $E$,
and the sketch is tightened so the next agent pass operates under a stronger contract.

The central claim is modest but operational:

\begin{quote}
A coding-agent workflow becomes inspectable when each counterexample is linked to the
sketch clause it supplements, the implementation or prompt path that handles it, and the
replay/compare check that would fail the tempting wrong patch.
\end{quote}

This claim is deliberately weaker than total correctness.  It is also stronger than a demo.  It
gives maintainers a local correctness theorem over the current corpus and a procedure for
growing that corpus when the theorem is falsified.

\paragraph{Contributions.}  This paper makes four contributions.

\begin{enumerate}[leftmargin=*]
  \item It defines counterexample-supplemented sketches as a repository-native artifact for
        agentic coding.
  \item It separates two hole-filling roles: Oracle~A, deterministic code the agent maintains,
        and Oracle~B, prompt/LLM/cassette behavior for sketch-declared narrative holes.
  \item It states a finite-corpus correctness theorem for replay/compare gates and proves it
        from the executable gate contract.
  \item It provides a runnable worked example, RosterSynth kiosk double-booking, where a
        replay-valid but semantically wrong repair becomes a counterexample that strengthens
        the sketch.
\end{enumerate}

% ---------------------------------------------------------------------------
\section{Lineage and departure from Sketch}
\label{sec:lineage}

Solar-Lezama's sketching work framed synthesis as a collaboration between programmer
insight and automated search: the programmer writes a partial program, leaving holes for a
synthesizer to fill~\cite{solar2008sketching,solar2013sketching}.  The point is not that the
human disappears.  The point is that the human supplies high-level structure while the
machine discharges low-level details.  The 2008 dissertation also emphasizes the central
challenge we inherit: once synthesis becomes search rather than derivation, correctness must
be established by an external validation procedure.  CEGIS addresses this by alternating
candidate generation with validation and feeding counterexamples back into the synthesis
loop.

Our setting changes the machinery but preserves the pressure.  The synthesizer is no longer
a SAT-backed Sketch compiler.  It is a coding agent inside a repository.  The holes are not
only integer constants or generator choices.  They include ordinary source edits, API choices,
error-shape conventions, and prompt instructions.  The validator is not a theorem prover.  It
is a repo-native gate: tests, replay checks, semantic comparison, fixture runs, and source
selectors.

\begin{center}
\begin{tabular}{p{0.22\linewidth}p{0.22\linewidth}p{0.22\linewidth}p{0.24\linewidth}}
\toprule
 & Sketch / CEGIS & Programming by example & This work \\
\midrule
Human input & Partial program with holes & Input/output examples, DSL & Sketch file, examples, counterexamples, known-code anchors \\
Synthesizer & Solver-backed search & DSL learner / witness functions & Coding agent editing repo files \\
Counterexample source & Validator / solver & User examples or generated examples & Gate failure promoted into corpus $E$ \\
Validation & Formal or bounded decision procedure & Consistency with examples & Replay + semantic compare + tests over $E$ \\
Guarantee & Solver-relative correctness & DSL/example consistency & Finite-corpus correctness under executable gate \\
\bottomrule
\end{tabular}
\end{center}

Programming-by-example systems such as FlashMeta and PROSE synthesize DSL programs from
examples using domain-specific deduction~\cite{polozov2015flashmeta,gulwani2017program}.
Agent benchmarks such as SWE-bench evaluate whether language models can resolve real
repository issues~\cite{jimenez2024swebench}.  Tests-as-prompts work studies tests as both
input and evaluation for LLM code generation~\cite{cui2025tests}.  Agentic synthesis
against counterexample-supplemented sketches sits between these traditions: it accepts that
agents edit ordinary code, but it refuses to make natural-language prompting the only
contract.

% ---------------------------------------------------------------------------
\section{The programming model}
\label{sec:model}

The method is easiest to understand as a repository contract.

\begin{definition}[Sketch]
A sketch $\mathcal{S}$ is a human-editable artifact that fixes the permitted strategy space
and names holes.  It may contain prose, tables, pseudo-code, type shapes, decision order,
and explicit abstention rules.  Unlike a chat prompt, it persists in the repository and is
edited when counterexamples reveal missing policy.
\end{definition}

\begin{definition}[Corpus]
The corpus $E = \{(x_i, y_i)\}_{i=1}^{n}$ is the finite set of observations the current gate
must satisfy.  Each input $x_i$ is a concrete scenario or fixture.  Each $y_i$ is the expected
semantic outcome: not merely the output shape, but the fields replay under-specifies.
\end{definition}

\begin{definition}[Counterexample-supplemented sketch]
A sketch is counterexample-supplemented when failures are promoted into $E$ and used to
revise $\mathcal{S}$.  The sketch is therefore not just an initial specification; it is the
stable surface on which the corpus teaches the next agent pass what was previously missing.
\end{definition}

\begin{definition}[Known-code anchor]
A known-code anchor is an existing source artifact that constrains the agent's implementation
style: an API, result type, parser convention, error shape, fixture format, test helper, or
reference implementation.  It prevents the agent from satisfying the example by inventing a
parallel local convention.
\end{definition}

The model distinguishes two implementation roles.

\paragraph{Oracle A: deterministic code.}
Oracle~A is source code maintained by the agent or developer.  It handles holes that are
encodable in ordinary control flow: selecting a duplicate booking, choosing a work date,
constructing a typed row, preserving an error convention.  Oracle~A should be boring.  It is
the part future maintainers can debug without consulting a model transcript.

\paragraph{Oracle B: narrative completion.}
Oracle~B is the LLM-mediated path for holes the sketch marks as narrative or non-encoding:
ambiguous notes, human descriptions, or policy text whose resolution is not yet worth
turning into deterministic code.  Oracle~B is still constrained by the sketch and checked by
the same gate.  In CI, a cassette can freeze the output so the gate remains reproducible.

\paragraph{Hybrid resolution.}
A hybrid resolver runs Oracle~A first.  If Oracle~A abstains or produces a row that replay
rejects, the resolver falls back to Oracle~B.  This keeps deterministic code on the hot path
while preserving an explicit route for narrative holes.

% ---------------------------------------------------------------------------
\section{Gate semantics}
\label{sec:gate-semantics}

The gate is the executable semantics of the corpus.  It has two layers.

\begin{definition}[Replay]
Replay applies a candidate output $r$ to an input state $x$ and checks whether the state
change closes the domain gap.  In RosterSynth, replay computes the effective coverage delta
after appends or booking modifications.  A row is replay-valid when the effective delta is
zero up to tolerance.
\end{definition}

\begin{definition}[Semantic compare]
Semantic compare checks fields that replay cannot infer: operation kind, selected booking,
work date, issue type, status mutation, or other policy-bearing columns.  A row is
compare-valid when it matches the golden expectation for that scenario.
\end{definition}

\begin{definition}[Gate pass]
For a strategy $H$ and corpus $E$, the gate passes iff for every $(x_i, y_i) \in E$, the
candidate row $r_i = H(x_i)$ is both replay-valid on $x_i$ and compare-valid against $y_i$.
\end{definition}

Replay and compare are intentionally redundant in some cases.  Replay prevents semantically
correct-looking rows that do not actually repair the state.  Compare prevents state-repairing
rows that violate policy.  The kiosk example needs both: appending $-40$ hours closes the
hours math, but the correct policy is to cancel the duplicate booking.

% ---------------------------------------------------------------------------
\section{Finite-corpus correctness}
\label{sec:correctness}

This section states the proof obligation honestly.  The method does not prove correctness
for all possible inputs.  It proves correctness over the current promoted corpus under the
repository's executable semantics.  That is the right theorem for a coding-agent workflow:
when a new failure appears, it becomes a new element of $E$, and the theorem's domain grows.

\begin{definition}[$E$-correctness]
A strategy $H$ is $E$-correct with respect to replay function $R$ and compare predicate $C$
when, for every $(x_i, y_i) \in E$, $R(x_i, H(x_i)) = \mathrm{true}$ and
$C(y_i, H(x_i)) = \mathrm{true}$.
\end{definition}

\begin{lemma}[Replay soundness relative to the implementation]
If the replay checker returns true for $(x, r)$, then $r$ satisfies the state-repair condition
encoded by the replay checker for $x$.
\end{lemma}

\begin{proof}
The replay checker is the executable definition of state repair.  In the RosterSynth example,
it evaluates appends and modifications against the observed roster state and returns true
only when the effective coverage delta is closed.  Therefore a true result entails exactly the
state-repair condition the checker implements.  This is implementation-relative soundness,
not an external proof that the checker captures every business rule.
\end{proof}

\begin{lemma}[Compare soundness relative to the corpus]
If semantic compare returns true for expected row $y$ and candidate row $r$, then $r$ agrees
with $y$ on every policy-bearing field implemented by the comparer.
\end{lemma}

\begin{proof}
The comparer enumerates the checked fields and fails on mismatch.  For modify rows in the
worked example, it checks the operation, issue type, booking id, and status fields.  Therefore
a true result means no checked field differed.  Again, the guarantee is relative to the
comparer's encoded field set.
\end{proof}

\begin{theorem}[Finite-corpus gate soundness]
If the gate passes for strategy $H$ on corpus $E$, then $H$ is $E$-correct with respect to the
replay and compare semantics encoded by the repository.
\end{theorem}

\begin{proof}
By definition, the gate iterates over every $(x_i, y_i) \in E$, computes $r_i = H(x_i)$, and
requires both replay and compare to pass.  By the replay lemma, each accepted $r_i$ satisfies
the state-repair condition encoded by replay for $x_i$.  By the compare lemma, each accepted
$r_i$ agrees with $y_i$ on every field encoded by semantic compare.  Therefore every corpus
case satisfies both predicates, which is exactly $E$-correctness.
\end{proof}

\paragraph{What the theorem does not say.}
It does not say that $H$ is correct on inputs outside $E$.  It does not say the replay checker
is a complete business specification.  It does not say the golden rows are perfect.  It says
something narrower and useful: the current artifact cannot be called done unless it satisfies
every promoted counterexample under the executable semantics the repository exposes.

\paragraph{Bounded observation hypothesis.}
The analogy to Sketch is the bounded observation hypothesis: a small set of carefully chosen
inputs can often constrain the desired implementation in practice~\cite{solar2008sketching}.
Here that hypothesis becomes operational rather than automatic.  The human and agent grow
$E$ when failures reveal missing policy.  The corpus is not assumed complete; it is the
current boundary of accountability.

% ---------------------------------------------------------------------------
\section{The synthesis loop}
\label{sec:loop}

The process is a CEGIS-shaped loop with an agent in the repair role.

\begin{enumerate}[leftmargin=*]
  \item Start with sketch $\mathcal{S}$, corpus $E$, and known-code anchors $K$.
  \item Ask the coding agent to implement or repair Oracle~A / Oracle~B while preserving
        $\mathcal{S}$ and $K$.
  \item Run the gate over $E$.
  \item If all cases pass, the artifact is $E$-correct.
  \item If a case fails, promote the failure into $E$, revise $\mathcal{S}$ so the missing rule
        is explicit, and repeat.
\end{enumerate}

The important step is promotion.  A failure that remains only in chat is not a counterexample;
it is a memory leak.  A failure becomes useful when it changes the corpus and the sketch so
future runs cannot ignore it.

\paragraph{Why agents fit the repair role.}
Agents are good at local synthesis under examples and constraints.  They are bad at
remembering invisible constraints across sessions.  Counterexample-supplemented sketches
turn hidden constraints into repository artifacts.  The agent can then do what it is good at:
apply local edits to code and prompts under explicit validation pressure.

% ---------------------------------------------------------------------------
\section{Worked example: RosterSynth kiosk double-booking}
\label{sec:worked-example}

RosterSynth is a synthetic roster-remediation reference implementation.  It is only a worked
example.  The process name is agentic synthesis against counterexample-supplemented
sketches.

\paragraph{Scenario.}
Employee E-1800 has 40 badge hours and 80 scheduled hours, so $\Delta = -40$.  The roster
contains twin active 40-hour bookings on the pay-window end date, ids 1801 and 1802.  The
correct repair is to cancel the higher duplicate booking, 1802, by emitting a modify row with
\texttt{fields.status=4}.

\paragraph{Tempting wrong patch.}
A naive implementation can append $-40$ hours.  Replay accepts this because the coverage
math closes.  Compare rejects it because the operation is wrong: the policy is duplicate
cancellation, not a compensating adjustment.

\paragraph{Promotion.}
The historical failure is:

\begin{quote}\small
\texttt{expected op modify, got append}; replay note: \texttt{all imbalanced rosters closed}.
\end{quote}

That failure promotes the kiosk scenario into $E$ and adds Op~2 to the sketch: twins on the
window end date with the same $(\mathrm{shiftKind}, \mathrm{hours})$ cancel the higher
\texttt{bookingId}, only if that cancellation alone closes $\Delta$.

\paragraph{Oracle A.}
The deterministic implementation groups active bookings by $(\mathrm{shiftKind},
\mathrm{hours})$, picks the maximum booking id in a duplicate group, checks whether
cancellation closes $\Delta$, and emits \texttt{op=modify} with \texttt{status=4}.

\paragraph{Oracle B.}
The prompt path states the same decision order: Op~2 before Op~1, higher booking id as the
default tie-break, and a required modify shape.  The cassette path lets CI run without a
live model.

\paragraph{Negative check.}
A wrong cassette that cancels booking 1801 still closes replay but fails compare:
\texttt{expected bookingId 1802, got 1801}.  This is the smallest demonstration of why
replay and compare must be separate layers.

% ---------------------------------------------------------------------------
\section{Reference artifact and reproducibility}
\label{sec:artifact}

The repository contains a full executable artifact for the worked example:

\begin{center}
\begin{tabular}{ll}
\toprule
Artifact & Role \\
\midrule
\texttt{examples/rostersynth-kiosk/source/docs/sketch.md} & Sketch clauses \\
\texttt{examples/rostersynth-kiosk/source/scenarios/*.json} & Corpus $E$ \\
\texttt{examples/rostersynth-kiosk/source/rostersynth/playbook.py} & Oracle~A \\
\texttt{examples/rostersynth-kiosk/source/rostersynth/oracle/prompt.py} & Oracle~B prompt \\
\texttt{examples/rostersynth-kiosk/source/cassettes/*.json} & Oracle~B cassette fixtures \\
\texttt{examples/rostersynth-kiosk/tests/test\_extracted\_rostersynth.py} & Executable checks \\
\texttt{build/rostersynth-kiosk-graph.json} & Extracted provenance graph \\
\bottomrule
\end{tabular}
\end{center}

Run:

\begin{verbatim}
python3 tools/extract_rostersynth_example.py --write --paper
python3 tools/extract_rostersynth_example.py --ce roster.kiosk_double_booking.v1
python3 -m unittest discover -s examples/rostersynth-kiosk/tests
\end{verbatim}

The test suite checks six claims: Oracle~A cancels booking 1802; wrong append passes replay
and fails compare; wrong cassette passes replay and fails compare; hybrid uses deterministic
for kiosk; the prompt carries decision order and payload; and the full corpus gates match the
recorded deterministic/hybrid/LLM-only cassette evidence.

% ---------------------------------------------------------------------------
\section{Discussion}
\label{sec:discussion}

\paragraph{Why not just tests?}
Tests can verify behavior, but they do not necessarily name the missing rule.  A
counterexample-supplemented sketch does both: the failing case joins $E$, and the sketch
records the rule the failure exposed.  The next agent pass sees the rule before editing.

\paragraph{Why not just prompts?}
Prompts are too easy to fork invisibly.  The sketch persists as a repo artifact.  Prompt text
is one implementation surface, not the only contract.

\paragraph{Why not formal synthesis?}
When policy is fully encodable, formal synthesis is stronger.  But many coding-agent tasks
mix encodable policy with conventions, prose, examples, and codebase-local style.  This
method targets that messy middle: stronger than ad-hoc prompting, weaker than universal
formal verification, and useful precisely because it fits ordinary repositories.

\paragraph{What counts as progress?}
The useful metric is not whether the agent produced code.  The useful metric is whether the
repo now contains a stronger sketch, a sharper counterexample corpus, and checks that reject
the previously tempting wrong patch.

% ---------------------------------------------------------------------------
\section{Limitations}
\label{sec:limitations}

\begin{enumerate}[leftmargin=*]
  \item \textbf{Finite corpus.}  The theorem is only over $E$.  New failures must be promoted.
  \item \textbf{Checker quality.}  Replay and compare are only as good as their encoded
        semantics.
  \item \textbf{Golden quality.}  A wrong golden row makes the gate enforce the wrong behavior.
  \item \textbf{Sketch quality.}  If the sketch hides a rule in prose nobody reads, the agent may
        still miss it.
  \item \textbf{Prompt drift.}  Oracle~B can drift when live model behavior changes; cassettes
        make CI reproducible but do not prove live stability.
  \item \textbf{Human judgment.}  The method surfaces obligations; it does not remove the need
        for domain review.
\end{enumerate}

% ---------------------------------------------------------------------------
\section{Conclusion}
\label{sec:conclusion}

Agentic synthesis against counterexample-supplemented sketches treats a coding agent as a
synthesizer inside a constraint environment.  The human writes the sketch, curates the corpus,
and promotes failures.  The agent edits code and prompts.  The gate supplies the finite-corpus
correctness boundary.  The result is not magic and not formal correctness over the universe.
It is a practical way to make agent-written code explainable, reusable, and progressively
harder to fool with the same wrong patch twice.

\bibliographystyle{plain}
\bibliography{references}

\end{document}
